\section{Package-Diagramm-Notation}


\subsection{Die Pakete}
\begin{description}
    \item[main:] Startpunkt des Systems; enthält \texttt{Main}.
    \item[docs:] Dokumente und Verwaltung; \texttt{Document}, \texttt{Drucksache}, \texttt{Rede}, \texttt{Protokoll}, \texttt{DocumentEditor}.
    \item[factory:] Erstellt Dokumente mit \texttt{DocumentFactory} und \texttt{DocumentData}.
    \item[logs:] Änderungs-Log mit \texttt{ChangeLogEntry} und \texttt{ChangeType}; keine Abhängigkeiten zu anderen Paketen.
    \item[storage:] Speichern und Laden; \texttt{DocumentStorage} und \texttt{JsonDocumentStorage} für JSON-Dateien.
    \item[users:] Benutzer und Berechtigungen; Interface \texttt{Berechtigung}.
    \item[users.human:] Menschliche Benutzer; \texttt{HumanUsers}, \texttt{Admin}, \texttt{DokumentBearbeiter}, \texttt{LogsBearbeiter}, \texttt{Abgeordneter}.
    \item[users.client:] Externe Anwendungen; \texttt{ExternalAPIClient}.
\end{description}

\subsection{Abhängigkeiten}
\begin{itemize}
    \item \texttt{docs} $\leftrightarrow$ \texttt{factory}: Dokumente erstellen.
    \item \texttt{docs} $\leftrightarrow$ \texttt{storage}: Dokumente speichern und laden.
    \item \texttt{docs} $\rightarrow$ \texttt{logs}: Änderungen protokollieren.
    \item \texttt{storage} $\rightarrow$ \texttt{logs}: Log-Einträge speichern und laden.
    \item \texttt{docs} $\rightarrow$ \texttt{users}: Benutzer und Berechtigungen verwenden.
    \item \texttt{users.human} $\rightarrow$ \texttt{users}: \texttt{HumanUsers} implementiert \texttt{Berechtigung}.
    \item \texttt{users.client} $\rightarrow$ \texttt{users}: \texttt{ExternalAPIClient} implementiert \texttt{Berechtigung}.
    \item \texttt{main} $\rightarrow$ \texttt{docs}, \texttt{factory}, \texttt{storage}, \texttt{users}: Startet und verbindet die Komponenten.
\end{itemize}

\subsection{Hinweise}
\begin{itemize}
    \item Abhängigkeiten zu \texttt{users.human} und \texttt{users.client} werden im Diagramm als ein Pfeil zu \texttt{users} gezeigt.
    \item Gegenseitige Abhängigkeiten zwischen \texttt{docs}, \texttt{factory} und \texttt{storage} sind für die Zusammenarbeit nötig.
    \item Die Setter der Dokumente sind package-private. Änderungen sollen nur über den \texttt{DocumentEditor} erfolgen, damit das Änderungs-Log immer aktuell bleibt.
    \item Java-Klassen, JSON-Bibliotheken und Tests sind nicht dargestellt.
    \item Erweiterbarkeit: Neue Dokumenttypen benötigen Änderungen in \texttt{docs}, \texttt{factory} und \texttt{storage}. Neue Benutzerrollen kommen in \texttt{users.human}.
\end{itemize}